% Work-MAT | 2. rocnik | hodina 17
\documentclass[11pt,a4paper]{article}
\input{../../sablony/pripravy/priprava.sty}
\usepackage{tikz}
\usetikzlibrary{arrows.meta}
\nastavPripravu{2. ročník}{Exponenciální rovnice}{17}{Exponenciální rovnice II: společný základ}
\begin{document}
\poznamka{Nácvik pravidel $a^{p}a^{q}=a^{p+q}$ a $(a^p)^q=a^{pq}$. Orientačně 6 min připomenutí, 12 min převody základů, 12 min obtížnější rovnice, 10 min samostatná práce, 5 min shrnutí.}
\blok{Nejdřív sjednotit základy}
Čísla $4,8,16,32$ jsou mocniny dvojky a $9,27,81$ mocniny trojky. Proto rovnice $4^x=8$ není vhodná pro přímé porovnání exponentů, dokud nepřepíšeme obě strany do stejného základu: $(2^2)^x=2^3$.
\begin{zadaniBox}[title={Věta - Pravidla pro mocniny}]
Pro $a>0$ a reálná $p,q$ platí $(a^p)^q=a^{pq}$, $a^pa^q=a^{p+q}$ a $\frac{a^p}{a^q}=a^{p-q}$. Mocniny s různými základy lze porovnat podle exponentů teprve po sjednocení na stejný základ $a\ne1$.
\end{zadaniBox}
\begin{prikladReseniBox}[title={Příklad 1}]
Vyřeš $9^{x-1}=27$.
\tcblower
\[\begin{aligned}(3^2)^{x-1}&=3^3\\3^{2x-2}&=3^3\\2x-2&=3\\x&=\frac{5}{2}.\end{aligned}\]
$K=\{\frac{5}{2}\}$.
\end{prikladReseniBox}
\begin{prikladReseniBox}[title={Příklad 2}]
Vyřeš $8^{x+1}=4^{2x-1}$.
\tcblower
\[\begin{aligned}(2^3)^{x+1}&=(2^2)^{2x-1}\\2^{3x+3}&=2^{4x-2}\\3x+3&=4x-2\\x&=5.\end{aligned}\]
Zkouška v exponentu: vlevo $3\cdot6=18$, vpravo $2\cdot9=18$; obě strany jsou $2^{18}$.
\end{prikladReseniBox}
\begin{prikladReseniBox}[title={Příklad 3}]
Vyřeš $\left(\frac{1}{4}\right)^{x-1}=8$.
\tcblower
\[\begin{aligned}(2^{-2})^{x-1}&=2^3\\2^{-2x+2}&=2^3\\-2x+2&=3\\x&=-\frac{1}{2}.\end{aligned}\]
$K=\{-\frac{1}{2}\}$.
\end{prikladReseniBox}
\begin{prikladReseniBox}[title={Příklad 4}]
Vyřeš rovnici
\[
\frac{16^{x-1}\cdot8^{x+1}}{2^x}=4^{2x+1}.
\]
Nejdříve vše převeď na základ 2 a přehledně sluč exponenty.
\tcblower
\[
\begin{aligned}
\frac{2^{4x-4}\cdot2^{3x+3}}{2^x}&=2^{4x+2},\\
2^{4x-4+3x+3-x}&=2^{4x+2},\\
6x-1&=4x+2,\\
2x&=3,\\
x&=\frac{3}{2}.
\end{aligned}
\]
$K=\left\{\frac{3}{2}\right\}$.
\end{prikladReseniBox}
\poznamka{Nezaměňuj $2^{x+2}$ za $2^x+2^2$. U převodu základu vždy nejdříve přepiš základ i exponent do závorky, například $(2^3)^{x-1}$.}
\begin{zadaniReseniBox}[title={Samostatná práce 1}]
\begin{enumerate}[label=\alph*)]
\item $25^{x-1}=125$.
\item $16^x=8^{x+2}$.
\item $\left(\frac{1}{9}\right)^x=27$.
\item $8^{x+1}=4^{2x-1}$.
\item $27^{x-2}=9^{x+1}$.
\item $\left(\frac{1}{4}\right)^{x-1}=2^{x+2}$.
\item $2^{x+2}\cdot4^{x-1}=32$.
\item $\frac{9^{x+1}}{3^{x-1}}=81$.
\end{enumerate}
\tcblower
\begin{enumerate}[label=\alph*)]
\item $5^{2x-2}=5^3$, tedy $x=\frac{5}{2}$.
\item $2^{4x}=2^{3x+6}$, tedy $x=6$.
\item $3^{-2x}=3^3$, tedy $x=-\frac{3}{2}$.
\item $2^{3x+3}=2^{4x-2}$, tedy $x=5$.
\item $3^{3x-6}=3^{2x+2}$, tedy $x=8$.
\item $2^{-2x+2}=2^{x+2}$, tedy $x=0$.
\item $2^{3x}=2^5$, tedy $x=\frac{5}{3}$.
\item $3^{2x+2-(x-1)}=3^{x+3}=3^4$, tedy $x=1$.
\end{enumerate}
\end{zadaniReseniBox}
\begin{zadaniReseniBox}[title={Rozšíření 1}]
Vyřeš $2^{x+2}\cdot4^{x-1}=32$.
\tcblower
$2^{x+2}\cdot2^{2x-2}=2^{3x}=2^5$, tedy $3x=5$ a $x=\frac{5}{3}$.
\end{zadaniReseniBox}
\clearpage
\begin{zadaniReseniBox}[title={Domácí úkol}]
\begin{enumerate}[label=\alph*)]
\item Vyřeš $27^x=9^{x+1}$.
\item Vyřeš $\left(\frac{1}{8}\right)^x=4^{x-1}$.
\item Vyřeš $8^{x-2}=16^{x+1}$.
\item Vyřeš $\frac{25^{x+1}}{5^{x-1}}=125$.
\item Vyřeš $\frac{2^{x+3}\cdot4^{x-1}}{8^x}=16$.
\item Proč nelze v rovnici $2^x=3^x$ hned porovnat exponenty?
\item Vysvětli význam závorek při převodu $8^{x-1}=(2^3)^{x-1}$.
\item Rozhodni, zda je $9^x=3^{2x}$ rovnicí s jediným kořenem.
\item Načrtni graf základní funkce $y=3^x$; vyznač čtyři body a asymptotu.
\item Před kreslením uprav $y=\frac{2^{x+3}}4-2$ na tvar se zřejmým posunem a sestroj jeho graf.
\end{enumerate}
\tcblower
\begin{enumerate}[label=\alph*)]
\item $3^{3x}=3^{2x+2}$, tedy $x=2$.
\item $2^{-3x}=2^{2x-2}$, tedy $x=\frac{2}{5}$.
\item $2^{3x-6}=2^{4x+4}$, tedy $x=-10$.
\item $5^{2x+2-(x-1)}=5^{x+3}=5^3$, tedy $x=0$.
\item $2^{x+3+2x-2-3x}=2^1=2$, nikoli 16; tedy $K=\varnothing$.
\item Základy jsou různé; nejprve musí být společný základ nebo jiný vhodný postup.
\item Mocnění mocniny dává $2^{3(x-1)}=2^{3x-3}$, ne $2^{3x-1}$.
\item Ne: jde o identitu, protože $9^x=(3^2)^x=3^{2x}$ pro všechna $x$.
\item Body $[-1;\frac{1}{3}]$, $[0;1]$, $[1;3]$, $[2;9]$; asymptota $y=0$.
\item $y=2^{x+1}-2$, takže posun o 1 doleva a o 2 dolů; asymptota $y=-2$. Body například $[-1;-1]$, $[0;0]$, $[1;2]$.
\end{enumerate}
\end{zadaniReseniBox}
\poznamka{Domácí úkol zhruba 20--30 minut: pět rovnic, tři otázky a dva grafy. U grafů žáci použijí dostatečně široké osy, podobně jako v hodině 13.}
\end{document}
